\section{From the three arithmetic inputs to an integral Case III seed}

% BEGIN S65A1
\paragraph{S65A1. Choosing the two tensor factors in the required order.}
Assume the three original inputs
$\mathrm{TauExAlt}$, $\mathrm{RhoExP}$ and
$\mathrm{RankFourAutomorphyOfTensor}$, with the respective standing
hypotheses $\mathrm{Af.Hypotheses}'$ and $\mathrm{PropEGP}$ needed
to apply the first two. Fix an arbitrary original Case III datum,
an acceptable tower $T_0$ with $\mathrm{Four}(T_0)$, the specified
complex comparison $\iota$ and the ambient coefficient datum.
The quadratic field $K'$ given by $\mathrm{Four}$ splits at $\ell$
and has discriminant prime to $\ell$. The rational image of $T_0.g_0$
lies outside $G_{K'}$. Otherwise its first root would be $1$ by
$\mathrm{QuadOn}$, contradicting the distinction between diagonal
characters $0$ and $1$ at the marked element.

Apply $\mathrm{TauExAlt}$ there. It supplies an exponent $m$, a
coefficient ring $O_0$, and a character $\kappa$ of $G_{K'}$ whose
positive powers have the required dihedral realizations over every
subsequent coefficient extension. Apply $\mathrm{RhoExP}$ with
this $O_0$ and with weight lower bound $2(9m)+3$. This gives one split
prime $w$, its auxiliary character $\theta$, a weight $k$, and the
ordinary rank-two representation $\rho$ with its specified automorphy,
open-subgroup irreducibility and $\mathrm{PR}$ condition. It also
gives an integral lattice $\rho_1$ and its slopes over one ring
$O_1$ extending both $O_0$ and the original coefficients. Its
congruence is $k-1\equiv e\bmod(\ell-1)$.

First specialize the dihedral promise over $O_1$ at power $1$.
This is the prerequisite for tensor automorphy to select a single
$1\le u\le9$. Now specialize the same promise at $u$ to obtain
$\tau$, its local signs, and $\kappa^u$ over $O_1$. The output of
tensor automorphy applies to this $\tau$ and this lattice since it
holds for every integral realization at the selected $u$. It gives
both the automorphic attachment away from the quadratic field
($\mathrm{AutOffS}$) and the additional assertion for every attached
cusp ($\mathrm{OffW}$). Put $M=mu$ (the product of $m$ and $u$).
Then $M\ge1$, by the positivity field of this dihedral realization,
and $2M+3\le k$. E65T2 supplies the correcting twist in the
original coefficient ring. These choices produce the original tensor
datum $d$, with dihedral weight $M$, and its additional $\mathrm{OffW}$
field. They precede the choice of every later tower.

Here is the ramification point in that construction. The auxiliary
character $\theta$ kills the group fixing $\mu_w$ and satisfies
$\theta^\ell=1$. Away from $w$, inertia acts trivially on
$\mu_w$, hence on $\theta$. The kernel is open since
$\mathbb Q(\mu_w)$ is finite. Identifying $\mu_\ell$ in the
coefficient field with the additive cyclic group of order $\ell$
does not change the kernel. Thus $\mathrm{Dies}(\ell,w,T.F)$,
where required later, puts the whole group $G_{T.F}$ in
$\ker\theta$.
% END S65A1

% BEGIN S65A2
\paragraph{S65A2. The characteristic-zero seed at every acceptable extension.}
Take the data of S65A1 and any acceptable $T\ge T_0$. Write
$P$ for the characteristic-zero realization of $Q_T$ and $R$ for
the untwisted signed tensor of $\rho_1$ and $\tau$. The rank-four
irreducibility result gives irreducibility of $P$ on
$G'=G_{T.L}\cap\ker\varepsilon$. Hence $R$ is irreducible on
$G_{T.L}$: any invariant subspace there would also be preserved
by the scalar twist of its restriction to $G'$. The field
$T.L=T.F\,T.K$ is abelian over $\mathbb Q$, as a compositum of two
abelian fields, and is totally complex. E65T3 excludes $K'\subset T.L$.
The local splitting and inertia argument for $\mathrm{GoodAt}$
in the preceding tensor argument therefore lets us apply
$\mathrm{AutOffS}$ and $\mathrm{OffW}$ to the very same restriction
$R|_{G_{T.L}}$. The first gives a regular algebraic cusp $\pi$ with
that attachment. The second supplies the original window condition
for every regular algebraic cusp with that attachment, hence for
this $\pi$. The character $\lambda=(\delta|_{G_{T.L}})^{-1}$ is
continuous and untwists $P$ to exactly that restriction.

Choose $j,U$ from the weight clause of the twist, so
$\delta^2=\varepsilon^j$ on the open subgroup $U$.
The integral weight calculation gives doubled weights
\[
 h_i=j+2a_i,\qquad
 (a_0,a_1,a_2,a_3)=(k-1+M,\ k-1-M,\ M,\ -M).
\]
Their successive half-gaps are $2M,k-1-2M,2M$, all at least two.
For the source parameter $a=-j$, $a+h_0=2(k-1+M)$ is even.
The rational character $\delta^{-2}\varepsilon^j$ has open
kernel containing $U$, so its image has finite order $b\ge1$.
Consequently $(\lambda^2\varepsilon^{-a})^b=1$ on all of
$G_{T.L}$, including every required inertia group. This gives
every field of $T.\mathrm{IsSeedAlt}(\iota,P)$.
% END S65A2

% BEGIN S65A3
\paragraph{S65A3. A single residual choice before all later integral seeds.}
Write $s$ for the image of $Q_{T_0}$ over the prescribed ambient
residue field. The reduction from $O_1$ kills its maximal ideal and,
on $O$, agrees with the original reduction by compatibility of the
ambient maps. On the initial frame with auxiliary set $\{w\}$,
E65S1 gives the determinant identity for this reduction on every
group element. Inertia away from $\ell,w$ kills $\theta$ by S65A1.
The integral pair calculation applied with this coefficient map
therefore gives the pair equations. E65S1 and E65S4 give the
residual diagonal, local characters, support and continuity. Thus
$s$ is residual on precisely the allowed initial frame and has
$N_M$ support.

For any later acceptable $T\ge T_0$ satisfying
$\mathrm{Dies}(\ell,w,T.F)$, use the same $O_1$ and take the
integral $Q_T$. Now $\theta$ is trivial on $G_{T.F}$, so the
integral determinant equation of the chosen rank-two lattice holds
on that entire group. The pair calculation with the identity
coefficient map applies to the empty auxiliary set. E65S4 gives
integral continuity, E65S1 gives the required residual conditions
over the residue field of $O_1$, and S65A2 gives the characteristic-zero
seed. Hence $Q_T$ satisfies $T.\mathrm{SeedPairAlt}$ over a finite
coefficient extension. Returning to the prescribed ambient residue
field, E65S3 gives compatibility with $s$ at every actual place,
E65S1 gives type $A$, and E65S2 gives the two root witnesses.
This proves the $\mathrm{SeedIVAlt}$ implication from the three
arithmetic inputs of S65A1 with their stated standing hypotheses.
The construction of those inputs is the remaining arithmetic task
in this reduction.
% END S65A3
